(1) 对$S=x^2+y^2-2z^2$有$\nabla^2S=2+2-4=0$，所以$\nabla^2\Phi=0$。离子受力$\boldsymbol F=-q\nabla\Phi$，于是$\ddot x+[q(U+V\cos\Omega t)/(mr_0^2)]x=0$、$\ddot y+[q(U+V\cos\Omega t)/(mr_0^2)]y=0$、$\ddot z-[2q(U+V\cos\Omega t)/(mr_0^2)]z=0$。三维不同方向的平均曲率满足电势Laplace约束，不能全部靠静电正定。
(2) 令$\tau=\Omega t/2$并用撇号表示$d/d\tau$，标准方程$u''+[a_u-2q_{M,u}\cos(2\tau)]u=0$。对$x,y$有$a_x=a_y=4qU/(mr_0^2\Omega^2)$、$q_{M,x}=q_{M,y}=-2qV/(mr_0^2\Omega^2)$；对$z$有$a_z=-8qU/(mr_0^2\Omega^2)$、$q_{M,z}=4qV/(mr_0^2\Omega^2)$。满足$a_x+a_y+a_z=0$、$q_{M,x}+q_{M,y}+q_{M,z}=0$，对应Laplace轨迹。
(3) 对$U=0$且$|q_{M,u}|\ll1$，将快速微运动与慢位移分离，周期平均给赝势$U_{\rm pseudo}=q^2|\boldsymbol E_{\rm rf}(\boldsymbol r)|^2/(4m\Omega^2)$，其中射频电场幅$\boldsymbol E_{\rm rf}=-(V/r_0^2)(x,y,-2z)$。因此$U_{\rm pseudo}=q^2V^2(x^2+y^2+4z^2)/(4m\Omega^2r_0^4)$，与$\tfrac12m(\omega_x^2x^2+\omega_y^2y^2+\omega_z^2z^2)$比较，得$\omega_x^2=\omega_y^2=q^2V^2/(2m^2\Omega^2r_0^4)$，$\omega_z^2=4\omega_x^2$，三个方向均正。
(4) 若$V=0$，静态四极势的Hessian迹为零，无法使三个方向曲率全为正；此即静电Earnshaw限制的局域体现。$V\ne0$时微运动从时间变化电场吸收并返回能量，平均赝势与$|\boldsymbol E_{\rm rf}|^2$成正，可实现动态三维约束。赝势近似须$|q_{M,u}|\ll1$、世俗频率远小于$\Omega$；在Mathieu不稳定舌、有限电极几何、离子碰撞和强微运动时必须采用完整Floquet稳定图。
